\section{Related Work}
\label{sec:related}

\subsection{Chromatin Interaction Prediction}

Chromatin interactions---most prominently captured by Hi-C~\cite{hic}---reflect the
three-dimensional folding of the genome and are mediated by architectural proteins
such as CTCF and cohesin~\cite{ctcf}.
Early computational approaches relied on epigenomic features: TargetFinder~\cite{targetfinder}
combined boosted trees with DNase-seq, histone mark ChIP-seq, CAGE, and DNA methylation
data, achieving strong cell-type-specific prediction but requiring expensive assays
unavailable for most conditions.

\subsection{Sequence-Based EPI Prediction}

The first deep learning method to predict EPI from raw DNA sequence was
SPEID~\cite{speid}, which processed one-hot encoded enhancer and promoter sequences
through stacked 1D convolutions followed by LSTM layers.
DeepC~\cite{deepc} extended this to megabase-scale chromatin contact prediction using
dilated convolutions and DNase-seq transfer learning.
EPI-Trans~\cite{epitrans} replaced convolutional backbones with a Transformer encoder,
reporting improved AUROC on multiple benchmarks.

\subsection{CNN and ResNet for Regulatory Genomics}

DeepSEA~\cite{deepsea} pioneered the application of CNNs to chromatin-profiling data,
learning sequence features predictive of transcription factor binding and chromatin
accessibility from 1 kb windows.
ResNet~\cite{resnet} architectures have since been adapted to genomics, providing residual
connections that facilitate training deep sequence models.

\subsection{RNN and LSTM Variants}

Classical BiLSTM~\cite{lstm} models have been widely applied in genomics to capture
directional sequence context.
xLSTM~\cite{xlstm} (Extended LSTM) introduces two advances over the classical model:
sLSTM (scalar memory with improved mixing) and mLSTM (matrix memory with a covariance
update rule, fully parallelisable), both stabilised via exponential gating.
The mLSTM variant is of particular interest for long regulatory sequences, as its matrix
memory $\mathbf{C}_t \in \mathbb{R}^{d_k \times d_v}$ retains richer contextual state
than the scalar LSTM cell.

\subsection{Transformer Variants for Long Sequences}

Standard self-attention has $O(L^2)$ complexity, prohibitive for raw 10 kbp sequences.
Three orthogonal solutions are incorporated in this work.

\textbf{Linear attention}~\cite{linear_transformer} re-expresses dot-product attention
via kernel feature maps, reducing complexity to $O(Ld^2)$ while preserving the global
receptive field.

\textbf{iTransformer}~\cite{itransformer} inverts the attention dimension: instead of
attending across positions, attention is applied across \emph{channels} (nucleotide
types), reducing complexity from $O(L^2 C)$ to $O(C^2 L)$, where $C = 5 \ll L$.

\textbf{BERT / DNABERT}~\cite{dnabert} adapts the bidirectional encoder representation
pre-trained with masked language modelling to 6-mer tokenised DNA.
The BERT encoder is the architectural backbone of all LLM-based encoders considered
in this work.

\subsection{Linear Recurrence and State Space Models}

\textbf{Mamba}~\cite{mamba} (Gu and Dao, 2023) introduces input-dependent (selective)
state space matrices, enabling dynamic filtering of informative content.
Mamba processes sequences in $O(L)$ time via hardware-efficient parallel associative
scans.

\textbf{RWKV}~\cite{rwkv} (Peng et al., 2023) reconciles Transformer parallelism
with RNN inference efficiency through a time-mixing mechanism based on element-wise
exponential decay, achieving $O(1)$ per-step cost at inference.

\textbf{HyenaDNA}~\cite{hyenadna} replaces attention with implicit long convolutions
(Hyena operators), supporting context lengths up to $10^6$ tokens at single-nucleotide
resolution.

\subsection{Self-Supervised Pre-training: MAE}

Masked Autoencoders (MAE)~\cite{mae} offer a complementary pre-training paradigm to
masked language modelling: a large fraction of input tokens (75\% in the original
vision work) is masked, and the encoder learns to reconstruct the missing content.
Unlike BERT-style objectives that predict discrete token identities, MAE targets
continuous-valued reconstruction, which may learn richer spatial representations.
Applied to DNA sequences, MAE pre-training can be used to initialise the Transformer
encoder with self-supervised sequence features before fine-tuning on labelled EPI
pairs.

\subsection{DNA Foundation Models}

DNABERT-2~\cite{dnabert2} replaces $k$-mer tokenisation with Byte Pair Encoding (BPE)
and scales pre-training to 135 species.
The Nucleotide Transformer~\cite{nucleotide_transformer} trains models up to 2.5B
parameters on 3,202 human genomes, yielding powerful low-data representations.

\subsection{Multi-Path Interaction Modelling}

A core design question in dual-sequence EPI models is how to combine enhancer and
promoter representations.
Simple concatenation~\cite{speid} ignores cross-sequence interaction, while element-wise
subtraction and multiplication capture complementary semantic differences and
co-activations.
Bilinear interaction layers provide the most expressive pairwise model.
We provide all of the above as configurable fusion options.
